\section{Toy Experiment Setup and Details}
\vspace{-0.3cm}
\label{sec:appx:toy exp setup and details}
In this section, we describe the setup for the toy experiment presented in the main paper. Specifically, the data distribution $p_0$ is a six-mode Gaussian mixture with covariance $0.3\mathbf{I}$, consisting of six equally weighted components uniformly arranged on a circle of radius 6. The reward function is defined as the sum of three Gaussian components centered at points evenly spaced on the circle: $r(x)=\sum_{i=1}^3\exp{(-0.5\|x-p_i \|^2})$, where $p_1=(6.5,0)$, $p_1=(-3.25, 3.25\sqrt{3})$, and $p_3=(-3.25, -3.25\sqrt{3})$. We trained a few-step score-based generative model using 4-layer MLP with hidden dimension 128. For sampling, we fixed the total NFE to 100 across all methods and visualized the results using 2,000 generated samples. We used step size 0.1 for MALA and 0.2 for pCNL. We deliberately used a smaller step size for MALA than for pCNL to better reflect the actual experimental settings used in the main experiments (Sec.~\ref{sec:experiments:setup}).